\cleardoublepage
\phantomsection
\addcontentsline{toc}{part}{Mixed Holomorphic--Topological Deligne Theory}
\chapter*{Book entrance: Mixed Holomorphic--Topological Deligne Theory}
\addcontentsline{toc}{chapter}{Book entrance: Mixed Holomorphic--Topological Deligne Theory}
\begin{summarybox}
Hamiltonian jets, shifted cotangent fields, universal branes, Duflo quantization, and the all-arity open--closed map.
\end{summarybox}
\part{The finite-jet laboratory}

\chapter{Moyal quantization in coordinates}
\section{The bidifferential formula}
For polynomials on $\C^2$, the Moyal product is
\begin{equation}\label{eq:vol2-moyal}
 f\star g
 =\sum_{n\ge0}\frac1{n!}\left(\frac{\hbarq}{2}\right)^n
 \sum_{r=0}^n(-1)^r\binom nr
 (\partial_1^{n-r}\partial_2^rf)
 (\partial_1^r\partial_2^{n-r}g).
\end{equation}
Every polynomial pair gives a finite sum.

\begin{theorem}[Quadratic exactness]
For polynomials $f,g$ of degree at most two,
\[
 [f,g]_\star=\hbarq\{f,g\}.
\]
The quadratic generators therefore define a Lie representation of the Rees-scaled $\mathfrak{sl}_2$.  Their associative algebra has an additional Casimir relation.
\end{theorem}
\begin{proof}
The commutator retains the odd powers of the Poisson bidifferential.  Derivatives of order three vanish on quadratics, so the first-order term is the entire commutator.
\end{proof}

\begin{proposition}[The oscillator quotient]\label{prop:vol2-oscillator-quotient}
Let $R_N=\C[\hbarq]/(\hbarq^{N+1})$ and let $W_N=R_N[q,p]$
have the Moyal product for $\{q,p\}=1$.
Let $U_N$ have generators $E,F,H$ with relations
\[
 [H,E]=2\hbarq E,\qquad [H,F]=-2\hbarq F,\qquad [E,F]=\hbarq H.
\]
Set $\Omega=EF+FE+H^2/2$.
The assignments
\[
 E\longmapsto e=q^2/2,\qquad
 F\longmapsto f=-p^2/2,\qquad
 H\longmapsto h=-qp
\]
induce an algebra isomorphism
\[
 U_N/(\Omega+3\hbarq^2/8)
 \xrightarrow{\sim} W_N^{(q,p)\mapsto(-q,-p)}.
\]
These compatible isomorphisms also hold after $\hbarq$-adic completion.
If linear Hamiltonians are retained, the generated associative algebra is the full Weyl algebra.
\end{proposition}
\begin{proof}
Quadratic exactness proves the three commutator relations.
The associative products are
\[
\begin{split}
 e\star f&=-q^2p^2/4-\hbarq qp/2-\hbarq^2/8,\\
 f\star e&=-q^2p^2/4+\hbarq qp/2-\hbarq^2/8,\\
 h\star h&=q^2p^2-\hbarq^2/4.
\end{split}
\]
Thus $\Omega$ maps to $-3\hbarq^2/8$.
Its commutator with each of $E,F,H$ vanishes by the defining relations.
Give these generators PBW degree one.
Give an even monomial $q^ap^b$ degree $(a+b)/2$.
The leading symbol of $\Omega+3\hbarq^2/8$ is
$C=2EF+H^2/2$.
It is a non-zero-divisor in $R_N[E,F,H]$, since $2C$ is monic in $H$.
Centrality and this property give
\[
 \operatorname{gr}\bigl(U_N/(\Omega+3\hbarq^2/8)\bigr)
 =R_N[E,F,H]/(C).
\]
Indeed, every element of the two-sided ideal is a multiple of the central generator.
Its leading symbol is the product of the two leading symbols.
The quotient has basis $E^aH^bF^c$ with $a,c\ge0$ and $b=0,1$.
Their leading symbols map to unit multiples of the distinct monomials
$q^{2a+b}p^{2c+b}$, precisely the even monomials.
Hence the associated-graded map is an isomorphism.
Lifting highest-degree terms proves surjectivity and comparing them proves injectivity.
Every element at fixed $N$ has finite polynomial degree.
The formulas commute with reduction in $N$, so their inverse limit gives the formal isomorphism.
Finally $q,p$ generate $W_N$ and topologically generate its formal limit.
\end{proof}

The primitive coproduct on $U_N$ does not descend to this quotient.
Modulo $\hbarq$, it sends the classical relation to
\[
 \Delta(C)=C\tensor1+1\tensor C
   +2E\tensor F+2F\tensor E+H\tensor H.
\]
The cross-term is nonzero in the tensor square of $\C[E,F,H]/(C)$.
For example, evaluation at $(E,F,H)=(1,0,0)$ in the first factor
and $(0,1,0)$ in the second gives $2$.
Thus the oscillator quotient is an associative algebra, with no coproduct inherited from these primitive generators.

\begin{calculation}[Cubic correction]
For $f=z_1^3$ and $g=z_2^3$,
\[
 f\star g-g\star f
 =9\hbarq z_1^2z_2^2+\frac32\hbarq^3.
\]
The first term quantizes the quartic Hamiltonian channel.  The constant is the first wheel-free ordering correction visible beyond the quadratic Lie algebra.
\end{calculation}
\begin{proof}
In \eqref{eq:vol2-moyal}, the odd orders are $n=1$ and $n=3$.  The first contributes $9\hbarq z_1^2z_2^2$.  The third contributes $2(\hbarq/2)^3(3!)=3\hbarq^3/2$.
\end{proof}

\section{Weyl generators}
With the convention $\{q,p\}=1$, the symbols satisfy
$[q,p]_\star=\hbarq$.
Their Weyl-ordered differential operators are
\[
 \widehat q=q,\qquad \widehat p=-\hbarq\partial_q,
 \qquad [\widehat q,\widehat p]=\hbarq.
\]
The abstract Weyl algebra is defined by this commutator relation.
Weyl ordering transports its multiplication to the Moyal product.
The differential-operator representation realizes these relations over
$\C[\hbarq]$ and its formal completion.
The finite quotient argument above uses the abstract $R_N$-algebra,
so it does not require a faithful differential-operator representation at nilpotent $\hbarq$.

\chapter{The quadratic invariant}
\section{The \texorpdfstring{$\mathfrak{sl}_2$}{sl2} invariant}
For the quadratic jet, identify the coordinate functions with $e,f,h$.  The invariant
\[
 q=2ef+\frac12h^2
\]
is $d_{\CE}$-closed.  Its formal Duflo image is $\Omega+\hbarq^2/2$, where $\Omega=ef+fe+h^2/2$ in the Rees enveloping algebra.  This is the degree-zero invariant comparison.  Under the oscillator quotient of Proposition~\ref{prop:vol2-oscillator-quotient}, its image is $\hbarq^2/8$.

\chapter{All-arity open--closed operations}
\section{Hochschild braces}
For Hochschild cochains $F\in C^p(A,A)$ and $G_i\in C^{q_i}(A,A)$, the brace
\[
 F\{G_1,\ldots,G_r\}
\]
is the signed sum over ordered insertions of the $G_i$ into the inputs of $F$.  The brace identities state that a nested insertion is the sum over all ways of distributing the inner insertions among the outer slots.

\begin{theorem}[Brace realization with a refinement]
For $A=A_d$, supply the homotopy-$\Etwo$ and mixed refinements of
Corollary~\ref{cor:vol2-open-closed-refinement}.
They identify the selected homotopy Gerstenhaber operations on
$\mathcal P_d$ with the relative Hochschild centre operations.
The chosen mixed refinement also identifies their action on the boundary algebra.
\end{theorem}
\begin{proof}
The relative Hochschild complex has its brace operations.
The supplied operadic equivalence preserves their coherent homotopies.
An equivalence of complexes alone would not imply this conclusion.
\end{proof}

\section{The first mixed identities}
A relative Hochschild $1$-cocycle $D$ of an algebra concentrated in degree zero
is an $R$-linear derivation:
\[
 D(ab)=D(a)b+aD(b).
\]
For two such cocycles $D,E$, their Gerstenhaber bracket acts by
\[
 [D,E](a)=D(E(a))-E(D(a)).
\]
For a graded algebra and homogeneous derivations, these formulas become
\[
\begin{split}
 D(ab)&=D(a)b+(-1)^{|D||a|}aD(b),\\
 [D,E](a)&=D(E(a))-(-1)^{|D||E|}E(D(a)).
\end{split}
\]
Higher Hochschild cochains act through multilinear evaluations and braces.
They need not act by derivations or carry a coproduct defining a diagonal action.
A compatible mixed refinement transports these operations, including their homotopies.

\begin{corollary}[Continuous open--closed structure with compatible jets]
Under the jet comparison and mixed refinement hypotheses of
Theorem~\ref{thm:vol2-compatible-jets}, the pair
$(\mathcal P_\infty,\widehat U_\hbar(\h))$
has the corresponding complete homotopy Swiss-cheese action.
Its closed color is equivalent to the continuous relative derived centre.
\end{corollary}
\begin{proof}
The supplied finite operations commute with the transition maps
and converge in the specified completion.
Their limit acts on the two limit objects.
The comparison theorem identifies the closed complex with continuous relative Hochschild cochains.
\end{proof}

\chapter{Finite Chan--Paton ranks}
\section{Rank one}
For $N=1$, the commuting equation is automatic.  The boundary coordinate algebra is $\C[x,y]$, and
\[
 J_1(f)=f(x,y).
\]
The necklace bracket reduces to the ordinary Poisson bracket.  Restriction to $f\in\mathfrak m^2$ gives the defining pointed-Hamiltonian representation.  The two linear Hamiltonians act by translations and recover the affine center-of-mass extension.

\section{Rank two and Cayley--Hamilton reduction}
For a $2\times2$ matrix $X$,
\[
 X^2-(\Tr X)X+(\det X)I=0.
\]
Hence every trace containing two consecutive powers of $X$ reduces to traces with $X$-degree at most one after adjoining $\Tr X$ and $\det X$.  The same statement holds for $Y$.  On the commuting locus, simultaneous diagonalization identifies the invariant algebra with the symmetric polynomials in two ordered pairs.

\begin{proposition}[First finite-rank trace relation]
For commuting $2\times2$ matrices,
\[
 p_{20}=p_{10}^2-2\det X,
 \qquad
 p_{02}=p_{01}^2-2\det Y,
\]
and
\[
 p_{11}^2-p_{20}p_{02}
 =-\bigl(x_1y_2-x_2y_1\bigr)^2.
\]
The stable generators remain independent until a relation forced by the rank is reached.
\end{proposition}
\begin{proof}
Write the joint eigenvalues as $(x_1,y_1)$ and $(x_2,y_2)$ and expand both sides.
\end{proof}

\section{Stable passage}
Adding the origin defines the rank transition.
Theorem~\ref{thm:vol2-stable-multisymmetric-ring} identifies the
stable graded algebra and its total-degree completion.
Proposition~\ref{prop:vol2-rank-poisson} retains the central rank
and identifies the stable Poisson sector of degree at least two.
A quotient map from a derived centre requires a representation
and trace comparison on the specified complexes; the stable
ring calculation alone does not construct that map.


\part{Global holomorphic symplectic descent}

\chapter{Fedosov descent}
\section{The Weyl algebroid}
All Hochschild cochains and Morita centres in this chapter are continuous
and relative to $R=\C\llbracket\hbarq\rrbracket$.
Let $(S,\omega)$ be a holomorphic symplectic surface.
Choose a formal Poisson tensor
$\pi_\hbar=\hbarq\omega^{-1}+O(\hbarq^2)$ and a Weyl algebroid
$\mathcal W_{S,\hbar}$, with specified holomorphic or Dolbeault
Fedosov resolutions and coherent overlap maps.
The formal Poisson class is part of these data.
It is not determined by the finite Hamiltonian Lie-cochain comparison.

\begin{theorem}[Global comparison with formality and descent data]
Suppose the resolutions carry the coherent complete
homotopy-$\Etwo$ comparison of
Theorem~\ref{thm:vol2-fedosov-marked-descent}, with
quasi-isomorphic associated graded maps. Then there is an equivalence
\[
 \Phi_S:
 \left(\PV_S\llbracket\hbarq\rrbracket,[\pi_\hbar,-]\right)
 \xrightarrow{\sim}
 C^\bullet_{\HH}(\mathcal W_{S,\hbar},\mathcal W_{S,\hbar}),
\]
and the corresponding equivalence on derived global sections.
A compatible gauge equivalence of the supplied resolutions
transports this comparison. Equivariance requires equivariant
comparison and descent data.
\end{theorem}
\begin{proof}
The complete filtered comparison is an equivalence on every
multiple intersection by Lemma~\ref{lem:vol2-complete-filtered-comparison}.
The homotopy limit of its descent diagram preserves that equivalence.
Apply derived global sections and the supplied gauge or equivariance
homotopies.
\end{proof}

\begin{corollary}[Global Morita centre with a continuous model]
Assume in addition that the chosen continuous Morita theory computes
the centre of $\Perf(\mathcal W_{S,\hbar})$ by derived global
sections of these relative Hochschild cochains. Then
\[
 Z_{\mathrm{Mor}}\bigl(\Perf(\mathcal W_{S,\hbar})\bigr)
 \simeq
 \R\Gamma\!
 \left(S,\PV_S\llbracket\hbarq\rrbracket,[\pi_\hbar,-]\right).
\]
This applies to a K3 surface when the same comparison data are supplied.
A Hamiltonian BF identification additionally requires a map from its
Chevalley observables to this symplectic polyvector complex.
\end{corollary}
\begin{proof}
Use the stated Morita model and the preceding derived global equivalence.
\end{proof}

\chapter{Factorization and descent}
\section{The algebraic mixed-HT realization}
Retain the complete homotopy-$\Etwo$ sheaf
\[
 \mathcal P_S=
 \left(\PV_S\llbracket\hbarq\rrbracket,[\pi_\hbar,-]\right).
\]
For each holomorphic open, its factorization realization uses two
framed topological directions. Suppose these realizations have the
cosheaf descent needed in the holomorphic variable.
They then define an algebraic mixed factorization theory
\[
 \Acal^{\mathrm{alg}}_{S,\mathrm{bulk}}
\quad\text{on}\quad
 \R^2_{\mathrm{top}}\times S.
\]
The supplied formality comparison realizes the corresponding
equivalence with the relative Hochschild centre.
An oriented realization requires rotation-equivariant structure.
Identifying $\Perf(\mathcal W_{S,\hbar})$ with a universal physical
boundary category requires the independent module, action and
generation comparisons of Chapter~\ref{ch:vol2-marked-reconstruction}.
Neither the factorization realization nor the scalar centre
calculation supplies them.

\section{BV realization}
The cyclic Lie sheaf of formal Hamiltonians has the shifted-cotangent action
\[
 S_{\mathrm{BF}}(A,B)
 =\int_{\R^2\times S}
 \left\langle B,(\dd+\dbar)A+\frac12[A,A]\right\rangle.
\]
Here $S$ is a complex symplectic surface, so $\mathbb R^2\times S$ has six real dimensions. A comparison of the classical observables with $\Acal^{\mathrm{alg}}_{S,\mathrm{bulk}}$ must specify a map from Hamiltonian Chevalley cochains to the symplectic polyvector complex. The finite Lie algebra identification does not supply it. Gauge fixing, extensions of graph weights and a solution of the quantum master equation would construct a perturbative factorization theory. Identifying its bulk--boundary action with the algebraic map additionally requires a comparison preserving support, units and collision homotopies. The finite jet zero modes carry the exact quantum measure constructed above. The three-dimensional half-space construction and its interacting comparison problem are treated separately in Chapter~\ref{ch:vol2-native-bulk-boundary}.

\part{Matrix branes and trace geometry}

\chapter{The derived commuting stack}
Let $X,Y\in\mathfrak{gl}_N$ and let $\psi\in\mathfrak{gl}_N[1]$.  The Koszul differential
\[
 QX=QY=0,
 \qquad
 Q\psi=[X,Y]
\]
resolves the commuting equation.  The simultaneous conjugation action gives the boundary cdga
\[
 A_N=C^*_{\CE}\!\left(
 \mathfrak{gl}_N,
 \Sym(\mathfrak{gl}_N\oplus\mathfrak{gl}_N\oplus\mathfrak{gl}_N[1])
 \right).
\]

Cyclic differentiation gives moment-preserving matrix derivations for
cyclic Hamiltonians. A polynomial Hamiltonian requires a cyclic lift,
and these lifts do not give a strict Lie map on all polynomials:
the cubic defect is \eqref{eq:vol2-cubic-trace-defect}.
Theorem~\ref{thm:vol2-all-arity-derived-action} constructs the
coherent action on the free resolution, and
Proposition~\ref{prop:vol2-action-bfv-extension} extends its matrix
realization through the gauge complex. Their higher relation
homotopies are part of the action. A strict continuous representation
of the completed enveloping algebra would require an additional
comparison in that same recipient.

\chapter{Necklace Hamiltonians}
Let
\[
 F=\C\angles{x,y},
 \qquad
 F_{\mathrm{cyc}}=F/[F,F].
\]
For a cyclic word $f$, define cyclic derivatives by cutting the word at every occurrence of the selected letter.  Put
\[
 \delta_f(x)=-\partial_y^{\mathrm{cyc}}f,
 \qquad
 \delta_f(y)=\partial_x^{\mathrm{cyc}}f.
\]

\begin{theorem}[Necklace Lie algebra]
The bracket
\[
 \{f,g\}_{\mathrm{neck}}
 =\left[
 \partial_x^{\mathrm{cyc}}f\,
 \partial_y^{\mathrm{cyc}}g-
 \partial_y^{\mathrm{cyc}}f\,
 \partial_x^{\mathrm{cyc}}g
 \right]_{\mathrm{cyc}}
\]
makes $F_{\mathrm{cyc}}/\C$ a Lie algebra.  The assignment $f\mapsto\delta_f$ is a Lie morphism.  It preserves the moment map:
\[
 \delta_f([x,y])=0.
\]
\end{theorem}
\begin{proof}
The cyclic Euler identity is
\[
 [x,\partial_x^{\mathrm{cyc}}f]+[y,\partial_y^{\mathrm{cyc}}f]=0.
\]
It gives the moment-map statement.  The canonical double Poisson bracket on the free algebra descends to the displayed necklace bracket; its Jacobi identity gives the first assertion.  Evaluation on generators gives $[\delta_f,\delta_g]=\delta_{\{f,g\}_{\mathrm{neck}}}$.
\end{proof}

\begin{theorem}[Trace moment map]
For
\[
 J_N(f)=\Tr f(X,Y),
\]
one has
\[
 \{J_N(f),J_N(g)\}_{\mathcal M_N}
 =J_N\!\left(\{f,g\}_{\mathrm{neck}}\right),
\]
where $\mathcal M_N=\mathfrak{gl}_N^2$ carries $\Tr(\dd X\wedge\dd Y)$.
\end{theorem}
\begin{proof}
The differential of a trace is the matrix evaluation of the cyclic derivative.  Contracting the two differentials with the inverse trace symplectic form gives the necklace formula.
\end{proof}

\chapter{The two-brane cone}
Separate center-of-mass and relative coordinates.  Put
\[
 u=x_1-x_2,
 \qquad
 v=y_1-y_2,
 \qquad
 E=\frac12u^2,
 \quad
 F=-\frac12v^2,
 \quad
 H=-uv.
\]
Then
\[
 H^2+4EF=0.
\]

\begin{proposition}[The $A_1$ Kirillov--Kostant cone]
With $\{u,v\}=2$,
\[
 \{H,E\}=4E,
 \qquad
 \{H,F\}=-4F,
 \qquad
 \{E,F\}=2H.
\]
After the evident rescaling this is the $\mathfrak{sl}_2$ Lie--Poisson bracket on the nilpotent cone.
\end{proposition}
\begin{proof}
The three identities follow by direct differentiation.  The quadratic equation is the determinant-zero equation for the traceless matrix
\[
 \begin{pmatrix}H&2E\\2F&-H\end{pmatrix}.
\]
\end{proof}

\part{Wheel corrections and exact quantum structure}

\chapter{The local centre with comparison data}
Assume a compatible jet comparison and a complete homotopy-$\Etwo$
refinement as in Theorem~\ref{thm:vol2-compatible-jets}.
Put
\[
 \mathcal A=\widehat U_\hbar(\h),\qquad
 \Acal_{\mathrm{bulk}}^{\mathrm{alg}}
 =\operatorname{Fact}_{\Etwo}(\mathcal P_\infty).
\]
The factorization envelope uses two framed topological directions.
An oriented refinement requires the corresponding rotation-equivariant structure.
The formal holomorphic coefficient remains at the Darboux point.
Let $\Ccal_\partial$ be the category of perfect complete $\mathcal A$-modules
in the chosen continuous $R$-linear Morita theory.
Assume its centre is computed by
$C^\bullet_{\HH,\mathrm{cont}}(\mathcal A/R,\mathcal A)$.

\begin{theorem}[Local derived-centre comparison]
These data give an equivalence of complete homotopy-$\Etwo$ factorization algebras
\[
 \Acal_{\mathrm{bulk}}^{\mathrm{alg}}
 \xrightarrow{\sim}
 Z^{\mathrm{der}}_{\mathrm{Mor},R}(\Ccal_\partial).
\]
At the formal insertion point the map is the selected refinement of $\Phi_\infty$.
Every continuous boundary module has the induced derived centre action.
\end{theorem}
\begin{proof}
The compatible jet theorem gives the complete homotopy-$\Etwo$ comparison.
The factorization envelope preserves this equivalence.
The assumed continuous Morita model identifies its target with the centre.
A centre element is a derived natural endomorphism of the identity functor,
so evaluation gives the action on each boundary module.
\end{proof}

\begin{corollary}[Global K3 comparison with descent data]
Let $S$ be a K3 surface with holomorphic symplectic form.
Suppose the local comparisons extend to a sheafwise equivalence with the
relative Hochschild centre of $\Perf(\mathcal W_{S,\hbar})$.
Suppose further that the selected cyclic and mixed refinements extend
and satisfy descent in the same completed category.
Then derived global sections give the corresponding global centre,
cyclic, and open--closed comparisons.
\end{corollary}
\begin{proof}
Derived global sections preserve equivalences of sheaves.
Apply them to the supplied sheafwise maps and their coherence homotopies.
\end{proof}

\chapter{Synthesis}
At each fixed Hamiltonian jet, the Rees enveloping algebra has relative
Hochschild complex equivalent to the Chevalley complex with differential
$\hbarq d_{\CE}$.
PBW antisymmetrization proves this over every $R_N$ and in the formal coefficient limit.
The Duflo theorem gives the separate multiplicative comparison on invariant functions.
Quadratic Moyal Hamiltonians generate the oscillator quotient of the enveloping algebra.
Their associative Casimir relation prevents the primitive coproduct from descending.

The Hamiltonian jet limit requires the compatible cochain systems of
Theorem~\ref{thm:vol2-compatible-jets}.
With its operadic refinement and continuous Morita model, the centre acts
on every supplied continuous boundary module $M$ through
\[
 \mathcal P_\infty
 \xrightarrow{\ \Phi_\infty\ }
 C^\bullet_{\HH,\mathrm{cont}}(\mathcal A/R,\mathcal A)
 \longrightarrow\RHom_{\mathcal A}(M,M).
\]
The last arrow evaluates a derived natural endomorphism at $M$.
A matrix-brane realization must construct such a module and compare its
centre action with the selected trace functions.
A commuting square with scalar representation invariants requires these
additional comparison maps.
Likewise, Fedosov descent concerns symplectic polyvectors and a Weyl algebroid.
Its relation to Hamiltonian Lie cochains requires the specified map between their carriers.

\part{Technical appendices to the preceding book}
\chapter{Bibliographic note}
The finite affine formality map and the separate symplectic globalization use Kontsevich formality and its Fedosov resolution \cite{KontsevichDQ,DolgushevCov,CalaqueDolgushevHalbout}.  The relative Hochschild model is the Cartan--Eilenberg comparison, with both Chevalley terms scaled by $\hbarq$.  Its finite-base resolution proof is Theorem~\ref{thm:vol2-relative-rees}.  Jet limits and operadic refinements require the additional data stated above.  The all-arity centre action is the Deligne--Tamarkin brace action \cite{McClureSmith,KontsevichSoibelmanDeligne}.  The global deformation-quantization category is the holomorphic Weyl algebroid.  The matrix trace and necklace formulas are the noncommutative symplectic moment-map construction.
